\documentclass[12pt, a4paper]{article}
\usepackage[utf8]{inputenc}
\usepackage[T1]{fontenc}
\usepackage[french]{babel}
\usepackage[margin=2.5cm]{geometry}
\usepackage{amsmath, amssymb}
\usepackage{listings}
\usepackage{xcolor}
\usepackage{tcolorbox}
\tcbuselibrary{theorems, skins, breakable}
\usepackage{booktabs}
\usepackage{fancyhdr}
\usepackage{hyperref}
\usepackage{tikz}
\usetikzlibrary{arrows.meta, positioning, shapes.geometric, chains}

% ── Couleurs ──────────────────────────────────────────────────────────
\definecolor{suva_blue}{RGB}{30, 100, 180}
\definecolor{code_bg}{RGB}{245, 247, 250}
\definecolor{code_kw}{RGB}{0, 80, 200}
\definecolor{code_str}{RGB}{180, 40, 40}
\definecolor{code_cmt}{RGB}{80, 130, 80}
\definecolor{warn_bg}{RGB}{255, 248, 220}
\definecolor{success_bg}{RGB}{230, 255, 235}
\definecolor{sol_green}{RGB}{0, 120, 80}
\definecolor{sol_bg}{RGB}{240, 255, 245}
\definecolor{block_bg}{RGB}{235, 240, 255}

% ── Styles listings ───────────────────────────────────────────────────
\lstdefinestyle{solidity}{
  language=C++,
  basicstyle=\ttfamily\small,
  keywordstyle=\color{code_kw}\bfseries,
  stringstyle=\color{code_str},
  commentstyle=\color{code_cmt}\itshape,
  backgroundcolor=\color{sol_bg},
  frame=single, rulecolor=\color{sol_green!40},
  numbers=left, numberstyle=\tiny\color{gray},
  breaklines=true, tabsize=4, showstringspaces=false,
  morekeywords={uint256, uint8, bytes, bytes32, bool, address, memory,
    storage, calldata, external, view, pure, returns, mapping, struct,
    event, emit, require, revert, modifier, payable, public, private,
    internal, virtual, override, assembly},
}
\lstdefinestyle{python}{
  language=Python,
  basicstyle=\ttfamily\small,
  keywordstyle=\color{code_kw}\bfseries,
  stringstyle=\color{code_str},
  commentstyle=\color{code_cmt}\itshape,
  backgroundcolor=\color{code_bg},
  frame=single, rulecolor=\color{gray!40},
  numbers=left, numberstyle=\tiny\color{gray},
  breaklines=true, tabsize=4, showstringspaces=false,
}
\lstdefinestyle{bash}{
  basicstyle=\ttfamily\small\color{white},
  backgroundcolor=\color{gray!80},
  frame=single, rulecolor=\color{gray!80},
  breaklines=true,
}

% ── En-tête / pied de page ────────────────────────────────────────────
\pagestyle{fancy}
\fancyhf{}
\fancyhead[L]{\color{suva_blue}\bfseries Blockchain \& Ethereum}
\fancyhead[R]{\color{gray}\small Leçon 11 --- SuVa Protocol}
\fancyfoot[L]{\footnotesize\color{gray}Mohamed Lamine MBENGUE}
\fancyfoot[C]{\thepage}
\fancyfoot[R]{\footnotesize\color{gray}portfolio-alamine.web.app}
\renewcommand{\footrulewidth}{0.3pt}

% ── Environnements tcolorbox ──────────────────────────────────────────
\newtcolorbox{definition}[1]{
  colback=suva_blue!5, colframe=suva_blue!60,
  title={\bfseries Définition : #1}, fonttitle=\color{white},
  breakable
}
\newtcolorbox{exemple}[1]{
  colback=success_bg, colframe=sol_green!60,
  title={\bfseries Exemple : #1}, fonttitle=\color{sol_green!80!black},
  breakable
}
\newtcolorbox{attention}{
  colback=warn_bg, colframe=orange!70,
  title={\bfseries Attention}, fonttitle=\color{orange!80!black},
  breakable
}
\newtcolorbox{pratique}[1]{
  colback=block_bg, colframe=suva_blue!80,
  title={\bfseries Exercice pratique : #1}, fonttitle=\color{white},
  breakable
}

% ═════════════════════════════════════════════════════════════════════
\begin{document}
% ═════════════════════════════════════════════════════════════════════

\begin{center}
  {\color{suva_blue}\rule{\textwidth}{2pt}}\\[0.5cm]
  {\Huge\bfseries\color{suva_blue} Blockchain \& Ethereum}\\[0.3cm]
  {\large\color{gray} De zéro à smart contract --- Leçon 11}\\[0.5cm]
  {\color{suva_blue}\rule{\textwidth}{0.5pt}}
\end{center}

\tableofcontents
\newpage

% ─────────────────────────────────────────────────────────────────────
\section{C'est quoi une blockchain ? (la base)}
% ─────────────────────────────────────────────────────────────────────

\begin{definition}{Blockchain}
Une blockchain (chaîne de blocs) est une \textbf{base de données partagée} entre des milliers
d'ordinateurs dans le monde. Personne ne la contrôle seul. Une fois qu'une donnée est écrite,
elle ne peut plus être modifiée ou supprimée.
\end{definition}

\subsection{L'analogie du cahier public}

Imagine un cahier posé en plein milieu d'une ville. N'importe qui peut lire ce qui est écrit.
N'importe qui peut écrire une nouvelle ligne. Mais personne ne peut effacer ou modifier ce qui
est déjà écrit. Et des milliers de personnes gardent une copie exacte de ce cahier chez eux.

C'est une blockchain.

\begin{center}
\begin{tikzpicture}[
  block/.style={draw, rounded corners, fill=block_bg, minimum width=3cm,
                minimum height=2cm, align=center, font=\small},
  arrow/.style={-Stealth, thick, suva_blue}
]
  \node[block] (b0) {\textbf{Bloc 0}\\\textit{(Genesis)}\\Hash: 0000...};
  \node[block, right=1.2cm of b0] (b1) {\textbf{Bloc 1}\\Tx: Alice→Bob\\Hash: a3f9...};
  \node[block, right=1.2cm of b1] (b2) {\textbf{Bloc 2}\\Tx: Bob→Carol\\Hash: 7c2e...};
  \node[block, right=1.2cm of b2] (b3) {\textbf{Bloc 3}\\Tx: ...\\Hash: d8b1...};

  \draw[arrow] (b0) -- (b1) node[midway, above, font=\tiny, gray]{contient hash de b0};
  \draw[arrow] (b1) -- (b2);
  \draw[arrow] (b2) -- (b3);
\end{tikzpicture}
\end{center}

\subsection{Pourquoi c'est sécurisé ?}

Chaque bloc contient :
\begin{itemize}
  \item Des transactions (qui a envoyé quoi à qui)
  \item Le \textbf{hash} (empreinte numérique) du bloc précédent
  \item Son propre hash
\end{itemize}

Si tu modifies un bloc, son hash change → le bloc suivant devient invalide → toute la chaîne
après est cassée. Et comme des milliers d'ordinateurs ont la même copie, ils verront
immédiatement que tu as triché.

\begin{attention}
La blockchain n'est pas magique. Elle est sécurisée parce qu'il est \textbf{économiquement
impossible} de contrôler plus de 50\% des ordinateurs du réseau en même temps.
\end{attention}

% ─────────────────────────────────────────────────────────────────────
\section{Ethereum --- La blockchain programmable}
% ─────────────────────────────────────────────────────────────────────

\begin{definition}{Ethereum}
Ethereum est une blockchain qui va plus loin que Bitcoin. En plus d'envoyer de l'argent,
on peut y exécuter des \textbf{programmes} appelés smart contracts. Ces programmes
tournent sur tous les ordinateurs du réseau en même temps.
\end{definition}

\subsection{Bitcoin vs Ethereum}

\begin{center}
\renewcommand{\arraystretch}{1.4}
\begin{tabular}{lll}
  \toprule
  \textbf{Aspect} & \textbf{Bitcoin} & \textbf{Ethereum} \\
  \midrule
  Usage principal & Envoyer de la monnaie & Monnaie + programmes \\
  Langage & Script limité & Solidity (Turing-complet) \\
  Monnaie native & BTC & ETH (Ether) \\
  Création & 2009 & 2015 \\
  Smart contracts & Non & Oui \\
  \bottomrule
\end{tabular}
\end{center}

\subsection{Les acteurs du réseau}

\begin{itemize}
  \item \textbf{Nœuds (nodes)} : ordinateurs qui stockent et valident la blockchain
  \item \textbf{Validateurs} : nœuds qui proposent et confirment les nouveaux blocs
  \item \textbf{Utilisateurs} : toi, moi --- on envoie des transactions
  \item \textbf{Smart contracts} : programmes déployés sur la blockchain
\end{itemize}

% ─────────────────────────────────────────────────────────────────────
\section{Les transactions}
% ─────────────────────────────────────────────────────────────────────

\begin{definition}{Transaction}
Une transaction est un message signé envoyé à la blockchain. Elle peut :
envoyer de l'ETH, appeler une fonction d'un smart contract, ou déployer un nouveau contrat.
\end{definition}

\subsection{Anatomie d'une transaction}

\begin{center}
\begin{tikzpicture}[
  field/.style={draw, fill=suva_blue!8, minimum width=8cm, minimum height=0.7cm,
                align=left, text width=7.5cm, font=\small}
]
  \node[field] (from) at (0,0) {\texttt{from}: 0xABC... (l'envoyeur)};
  \node[field, below=0.05cm of from] (to) {\texttt{to}: 0xDEF... (le destinataire)};
  \node[field, below=0.05cm of to] (value) {\texttt{value}: 0.1 ETH (montant)};
  \node[field, below=0.05cm of value] (data) {\texttt{data}: 0x... (données / appel de fonction)};
  \node[field, below=0.05cm of data] (gas) {\texttt{gasLimit}: 21000 (max gaz à utiliser)};
  \node[field, below=0.05cm of gas] (sig) {\texttt{signature}: r, s, v (prouve que tu es bien l'envoyeur)};
\end{tikzpicture}
\end{center}

\subsection{Le gas --- payer pour calculer}

\begin{definition}{Gas}
Le gas est l'unité qui mesure le travail de calcul sur Ethereum. Chaque opération coûte
un certain nombre de gas. Tu paies ce gas en ETH. Si ton programme est trop long ou
coûteux, la transaction échoue mais tu paies quand même le gas consommé.
\end{definition}

\begin{center}
\renewcommand{\arraystretch}{1.4}
\begin{tabular}{lrr}
  \toprule
  \textbf{Opération} & \textbf{Gas} & \textbf{Commentaire} \\
  \midrule
  Transfert ETH simple & 21\,000 & Le minimum \\
  Écrire dans le storage & 20\,000 & Très cher \\
  Lire depuis le storage & 2\,100 & Moins cher \\
  Addition (ADD) & 3 & Quasi gratuit \\
  Vérification ECDSA & 3\,000 & Signature actuelle d'Ethereum \\
  Vérification Dilithium & $\sim$6\,600\,000 & Post-quantique, cher \\
  Vérification Falcon & $\sim$350\,000 & Post-quantique, acceptable \\
  \bottomrule
\end{tabular}
\end{center}

\begin{exemple}{Calcul du coût réel}
Gas utilisé : 21\,000\\
Gas price : 10 gwei (1 gwei = $10^{-9}$ ETH)\\
Coût : $21\,000 \times 10 \times 10^{-9}$ ETH $= 0.00021$ ETH $\approx$ \$0.50
\end{exemple}

% ─────────────────────────────────────────────────────────────────────
\section{Adresses et clés sur Ethereum}
% ─────────────────────────────────────────────────────────────────────

\subsection{Comment fonctionne une adresse Ethereum}

\begin{center}
\begin{tikzpicture}[
  box/.style={draw, rounded corners, fill=block_bg, align=center, font=\small,
              minimum width=3.5cm, minimum height=1cm},
  arrow/.style={-Stealth, thick}
]
  \node[box] (priv) {Clé privée\\(256 bits aléatoires)};
  \node[box, right=1.5cm of priv] (pub) {Clé publique\\(courbe elliptique secp256k1)};
  \node[box, right=1.5cm of pub] (addr) {Adresse\\0x + 20 derniers octets\\du Keccak(clé pub)};

  \draw[arrow] (priv) -- (pub) node[midway, above, font=\tiny]{ECDSA};
  \draw[arrow] (pub) -- (addr) node[midway, above, font=\tiny]{Keccak-256};
\end{tikzpicture}
\end{center}

\begin{attention}
La clé privée est \textbf{secrète}. Celui qui la possède contrôle tous les fonds.
On ne peut pas retrouver la clé privée depuis l'adresse publique (sens unique).\\[0.2cm]
\textbf{Problème quantique} : un ordinateur quantique peut retrouver la clé privée depuis
la clé publique. C'est exactement ce que SuVa cherche à résoudre.
\end{attention}

\subsection{Signer une transaction}

Quand tu envoies une transaction, tu la signes avec ta clé privée. Ethereum vérifie
que la signature correspond à l'adresse \texttt{from}. C'est l'algorithme ECDSA.

\begin{center}
\begin{tikzpicture}[
  box/.style={draw, rounded corners, fill=block_bg, align=center, font=\small,
              minimum width=3cm, minimum height=1cm},
  arrow/.style={-Stealth, thick, suva_blue}
]
  \node[box] (tx) {Transaction\\(non signée)};
  \node[box, right=1.5cm of tx] (sign) {Signature\\ECDSA(clé privée, tx)};
  \node[box, right=1.5cm of sign] (net) {Réseau Ethereum\\vérifie avec clé publique};

  \draw[arrow] (tx) -- (sign);
  \draw[arrow] (sign) -- (net);
\end{tikzpicture}
\end{center}

% ─────────────────────────────────────────────────────────────────────
\section{Smart Contracts}
% ─────────────────────────────────────────────────────────────────────

\begin{definition}{Smart Contract}
Un smart contract est un programme stocké sur la blockchain. Il s'exécute automatiquement
quand on l'appelle. Ses règles sont publiques, son exécution est garantie. Personne
--- même le créateur --- ne peut le modifier après déploiement (sauf s'il est upgradeable).
\end{definition}

\subsection{Analogie : distributeur automatique}

Un distributeur de billets :
\begin{itemize}
  \item Tu insères ta carte (= tu envoies une transaction)
  \item Tu entres ton code (= tu signes)
  \item La machine vérifie (= le contrat vérifie)
  \item Elle te donne l'argent (= le contrat exécute)
\end{itemize}
Pas besoin de banquier humain. Le programme fait tout.

\subsection{Premier smart contract en Solidity}

\begin{lstlisting}[style=solidity, caption={Coffre-fort simple --- Vault.sol}]
// SPDX-License-Identifier: MIT
pragma solidity ^0.8.25;

contract SimpleVault {
    // owner = adresse qui a deploye le contrat
    address public owner;

    // Les fonds sont stockes dans le contrat lui-meme
    constructor() {
        owner = msg.sender; // msg.sender = celui qui deploie
    }

    // Deposer de l'ETH dans le coffre
    function deposit() external payable {
        // payable = cette fonction peut recevoir de l'ETH
        // msg.value = montant envoye
    }

    // Retirer ses fonds
    function withdraw(uint256 amount) external {
        require(msg.sender == owner, "Pas le proprietaire");
        require(address(this).balance >= amount, "Fonds insuffisants");

        (bool success, ) = msg.sender.call{value: amount}("");
        require(success, "Transfert echoue");
    }

    // Voir le solde du coffre
    function getBalance() external view returns (uint256) {
        return address(this).balance;
    }
}
\end{lstlisting}

\subsection{Les concepts Solidity essentiels}

\begin{center}
\renewcommand{\arraystretch}{1.4}
\begin{tabular}{lp{9cm}}
  \toprule
  \textbf{Concept} & \textbf{Explication} \\
  \midrule
  \texttt{msg.sender} & Adresse qui appelle la fonction \\
  \texttt{msg.value} & Montant ETH envoyé avec la transaction \\
  \texttt{storage} & Variables stockées sur la blockchain (cher) \\
  \texttt{memory} & Variables temporaires dans la transaction (moins cher) \\
  \texttt{require()} & Si faux → la transaction échoue et est annulée \\
  \texttt{event} & Log enregistré sur la blockchain (pas dans le storage) \\
  \texttt{mapping} & Dictionnaire clé → valeur (ex: adresse → solde) \\
  \texttt{payable} & Fonction qui peut recevoir de l'ETH \\
  \bottomrule
\end{tabular}
\end{center}

% ─────────────────────────────────────────────────────────────────────
\section{ERC-20 --- Les tokens}
% ─────────────────────────────────────────────────────────────────────

\begin{definition}{Token ERC-20}
Un token ERC-20 est un smart contract qui simule une monnaie. USDT, USDC, et même
des tokens SuVa (qUSDT) sont des ERC-20. Le standard définit les fonctions minimales
que tout token doit implémenter.
\end{definition}

\begin{lstlisting}[style=solidity, caption={Interface ERC-20 minimale}]
interface IERC20 {
    function totalSupply() external view returns (uint256);
    function balanceOf(address account) external view returns (uint256);
    function transfer(address to, uint256 amount) external returns (bool);
    function approve(address spender, uint256 amount) external returns (bool);
    function transferFrom(address from, address to, uint256 amount)
        external returns (bool);
}
\end{lstlisting}

\subsection{Lien avec SuVa}

\begin{center}
\begin{tikzpicture}[
  box/.style={draw, rounded corners, fill=block_bg, align=center, font=\small,
              minimum width=3cm, minimum height=1cm},
  arrow/.style={-Stealth, thick, suva_blue}
]
  \node[box] (usdt) {USDT\\\textit{(ERC-20 standard)}};
  \node[box, right=2cm of usdt] (vault) {QuantumVault.sol\\\textit{(smart contract SuVa)}};
  \node[box, right=2cm of vault] (qusdt) {qUSDT\\\textit{(ERC-20 SuVa)}};

  \draw[arrow] (usdt) -- (vault) node[midway, above, font=\tiny]{deposit()};
  \draw[arrow] (vault) -- (qusdt) node[midway, above, font=\tiny]{mint()};
  \draw[arrow] (qusdt) to[bend right=30] node[midway, below, font=\tiny]{burn() + withdraw()} (usdt);
\end{tikzpicture}
\end{center}

% ─────────────────────────────────────────────────────────────────────
\section{Foundry --- Tester des smart contracts}
% ─────────────────────────────────────────────────────────────────────

\begin{definition}{Foundry}
Foundry est le framework de développement Solidity utilisé dans ce projet. Il permet de
compiler, tester, et déployer des contrats. Les tests sont écrits en Solidity (pas en
JavaScript comme avec Hardhat).
\end{definition}

\subsection{Structure d'un projet Foundry}

\begin{lstlisting}[style=bash]
mon_projet/
|-- src/          # Contrats Solidity
|-- test/         # Tests Solidity
|-- lib/          # Dependances (forge-std, openzeppelin...)
|-- foundry.toml  # Configuration
\end{lstlisting}

\subsection{Écrire un test Foundry}

\begin{lstlisting}[style=solidity, caption={Test du SimpleVault}]
// SPDX-License-Identifier: MIT
pragma solidity ^0.8.25;

import "forge-std/Test.sol";
import "../src/SimpleVault.sol";

contract SimpleVaultTest is Test {

    SimpleVault vault;
    address alice = address(0x1);

    function setUp() public {
        // Donner de l'ETH a alice (dans le contexte de test)
        vm.deal(alice, 10 ether);

        // Deployer le vault avec alice comme proprietaire
        vm.prank(alice);
        vault = new SimpleVault();
    }

    function test_deposit() public {
        // Alice depose 1 ETH
        vm.prank(alice);
        vault.deposit{value: 1 ether}();

        // Verifier que le solde est correct
        assertEq(vault.getBalance(), 1 ether);
    }

    function test_withdraw() public {
        // Setup : alice depose
        vm.startPrank(alice);
        vault.deposit{value: 2 ether}();

        // Alice retire 1 ETH
        vault.withdraw(1 ether);
        vm.stopPrank();

        assertEq(vault.getBalance(), 1 ether);
    }

    function test_withdraw_unauthorized() public {
        address attacker = address(0x2);
        vm.prank(attacker);

        // Doit echouer car attacker n'est pas owner
        vm.expectRevert("Pas le proprietaire");
        vault.withdraw(1 ether);
    }
}
\end{lstlisting}

\subsection{Commandes Foundry essentielles}

\begin{lstlisting}[style=bash]
forge build                    # Compiler les contrats
forge test                     # Lancer tous les tests
forge test -vv                 # Mode verbose (voir les logs)
forge test --match-test test_deposit  # Tester une seule fonction
forge test --gas-report        # Rapport de consommation gas
anvil                          # Lancer un node Ethereum local
forge script deploy.s.sol      # Executer un script de deploiement
\end{lstlisting}

% ─────────────────────────────────────────────────────────────────────
\section{Exercices pratiques}
% ─────────────────────────────────────────────────────────────────────

\begin{pratique}{B1 --- Comprendre les transactions}
Lis une transaction réelle sur Ethereum:
\begin{enumerate}
  \item Va sur \texttt{etherscan.io}
  \item Cherche n'importe quelle adresse (ex: Vitalik: \texttt{0xd8dA6BF26964aF9D7eEd9e03E53415D37aA96045})
  \item Clique sur une transaction
  \item Identifie : \texttt{from}, \texttt{to}, \texttt{value}, \texttt{gas used}, \texttt{gas price}
  \item Calcule le coût en ETH et en USD
\end{enumerate}
\end{pratique}

\begin{pratique}{B2 --- Déployer SimpleVault avec Foundry}
\begin{lstlisting}[style=bash]
# 1. Creer le projet
mkdir my_vault && cd my_vault
forge init

# 2. Copier SimpleVault.sol dans src/
# 3. Copier le test dans test/

# 4. Compiler et tester
forge build
forge test -vv

# 5. Lancer anvil (node local)
anvil

# 6. Deployer sur le node local
forge script script/Deploy.s.sol --rpc-url http://localhost:8545 --broadcast
\end{lstlisting}
\end{pratique}

\begin{pratique}{B3 --- Mesurer le gas}
\begin{lstlisting}[style=bash]
# Voir exactement combien de gas chaque fonction consomme
forge test --gas-report

# Tu verras quelque chose comme:
# | Function   | min   | avg   | max   | # calls |
# | deposit    | 21000 | 21000 | 21000 | 3       |
# | withdraw   | 25000 | 25000 | 25000 | 2       |
\end{lstlisting}
Essaie d'ajouter une ligne \texttt{emit Deposit(msg.sender, msg.value);} dans
\texttt{deposit()} et vois comment le gas change.
\end{pratique}

\begin{pratique}{B4 --- Lire le code ETHDILITHIUM (Phase 1 SuVa)}
\begin{lstlisting}[style=bash]
cd falcon_dilithium/ETHDILITHIUM-main

# 1. Lancer les tests Solidity
forge test -vv

# 2. Voir le rapport gas
forge test --gas-report

# 3. Lancer les tests Python
cd python-ref
make test_signer
\end{lstlisting}
Questions à noter pendant la lecture :
\begin{itemize}
  \item Combien de gas pour \texttt{ZKNOX\_ethdilithium} sur ta machine ?
  \item Est-ce que tous les tests passent ?
  \item Quelle est la différence de gas entre Dilithium et ETHDilithium ?
\end{itemize}
\end{pratique}

% ─────────────────────────────────────────────────────────────────────
\section{Le lien avec SuVa}
% ─────────────────────────────────────────────────────────────────────

\subsection{Pourquoi la blockchain pose un problème de sécurité quantique}

Ethereum utilise ECDSA pour signer les transactions. La sécurité d'ECDSA repose sur la
difficulté de résoudre le \textit{problème du logarithme discret sur courbe elliptique}.
Un ordinateur quantique avec l'algorithme de Shor peut résoudre ce problème en temps
polynomial → toutes les signatures ECDSA peuvent être forgées.

\subsection{Ce que SuVa construit concrètement}

\begin{center}
\renewcommand{\arraystretch}{1.4}
\begin{tabular}{lll}
  \toprule
  \textbf{Composant} & \textbf{Fichier} & \textbf{Phase} \\
  \midrule
  Vérificateur Dilithium & \texttt{ZKNOX\_ethdilithium.sol} & Fait \\
  Signing Python & \texttt{Falcon\_dilithium\_eth.py} & Fait \\
  Vérificateur Falcon & À intégrer depuis ETHFALCON & Phase 2 \\
  Vault contract (qUSDT) & \texttt{QuantumVault.sol} (à créer) & Phase 3 \\
  Wallet ERC-4337 & \texttt{SuVaWallet.sol} (à créer) & Phase 4 \\
  \bottomrule
\end{tabular}
\end{center}

\subsection{Résumé de ce que tu dois savoir faire}

\begin{tcolorbox}[colback=suva_blue!5, colframe=suva_blue]
\begin{enumerate}
  \item \textbf{Lire} un smart contract Solidity et comprendre ce qu'il fait
  \item \textbf{Écrire} des fonctions \texttt{deposit()}, \texttt{withdraw()}, \texttt{verify()}
  \item \textbf{Tester} avec Foundry (\texttt{forge test})
  \item \textbf{Mesurer} le gas de chaque fonction
  \item \textbf{Déployer} sur un testnet (Sepolia)
  \item \textbf{Lier} le code Python (signing) avec le Solidity (verification)
\end{enumerate}
\end{tcolorbox}

% ─────────────────────────────────────────────────────────────────────
\section*{Glossaire rapide}
% ─────────────────────────────────────────────────────────────────────

\begin{center}
\renewcommand{\arraystretch}{1.4}
\begin{tabular}{lp{10cm}}
  \toprule
  \textbf{Terme} & \textbf{Définition simple} \\
  \midrule
  Block & Paquet de transactions validées ensemble \\
  Hash & Empreinte numérique unique d'une donnée \\
  Gas & Unité de mesure du calcul sur Ethereum \\
  Gwei & $10^{-9}$ ETH, unité du gas price \\
  EOA & Externally Owned Account --- un wallet classique \\
  Smart Contract & Programme déployé sur la blockchain \\
  ABI & Application Binary Interface --- décrit les fonctions d'un contrat \\
  Testnet & Réseau de test (Sepolia) --- ETH gratuit, pas de valeur réelle \\
  Mainnet & Réseau principal --- ETH a une valeur réelle \\
  Etherscan & Explorateur de la blockchain Ethereum \\
  Foundry & Framework de dev Solidity (forge, cast, anvil) \\
  Anvil & Node Ethereum local pour tester \\
  ERC-20 & Standard pour les tokens fongibles \\
  ERC-4337 & Standard pour les smart contract wallets \\
  \bottomrule
\end{tabular}
\end{center}

\vspace{1cm}

% ─────────────────────────────────────────────────────────────────────
\begin{tcolorbox}[colback=suva_blue!5, colframe=suva_blue!40,
  left=6pt, right=6pt, top=4pt, bottom=4pt]
\small\begin{tabular}{ll}
  \textbf{Auteur} & Mohamed Lamine MBENGUE \\
  \textbf{Spécialités} & Sécurité · DevOps · Dev Full Stack Web \& Mobile \\
  \textbf{Contact} & +221 78 178 44 36 · mbenguemohamedlamine2022@gmail.com \\
  \textbf{Portfolio} & \href{https://portfolio-alamine.web.app}{\texttt{portfolio-alamine.web.app}} \\
\end{tabular}
\end{tcolorbox}

\end{document}
